\scanpage{63}
\mspage{54}
% archive sheet number: 63
Si $\mathrel{|\circ|}$ dans $\unsure{M|\mathcal{M}}$ peut être défini par $\overset{\circ}{\ll}$, il n'est pas clair qu'il en soit de $=$ dans $M$, mais ce sera le cas si $M'$ est $\overset{\circ}{\ll}$-fermé (mais \inserted{bien} souvent, $M'$ ne le sera pas\ldots). Par contre, souvent on aura $M'\supset\unsure{L|\mathcal{L}}$, de sorte que $L$ est aussi l'un des lieux de $M'$. Si alors $M$ est \unsure{str. dis|str. disc}, $M'$ aussi (axiomatique $\mathrm{Mg}$ \unsure{disc|dis}, page 23) ; \deleted{\unsure{et}} si $M$ est \unsure{quasi-ensembliste}, $M'$ aussi. Enfin si $M$ \unsure{ensembliste}, $M'$ aussi\ldots

\figure{\begin{tikzpicture}[x=1cm,y=1cm]
\draw[line width=.5pt] (0,0) .. controls (.18,-.16) and (.4,-.48) .. (.53,-.8);
\end{tikzpicture}}

Je suis mûr maintenant pour étudier la notion d'homomorphismes \inserted{(\unsure{restreints})} entre atlas, donc des hom. qui respectent, non seulement $\overset{\circ}{\ll}$ et $\mathrel{|\circ|}$, mais aussi $\leq$. On les appellera provisoirement ``homomorphismes''.

\deleted{\uline{Définition} Soit}

Soient donc
\[
\begin{aligned}
&(\unsure{\mathscr{F}|\mathcal{F}|F},M,\triangleleft,\overset{\circ}{\ll},\mathrel{|\circ|})\\
&(\mathscr{F}',M',\triangleleft,\overset{\circ}{\ll},\mathrel{|\circ|})
\end{aligned}
\]
deux atlas, on veut définir les applications
\[
f:\mathscr{F}\longrightarrow\mathscr{F}'
\]
qualifiées d'\uline{homomorphismes}.
